\documentclass[11pt]{article}

\usepackage[spanish]{babel}
\usepackage[margin=1in]{geometry}
\usepackage{amsmath, amssymb, amsfonts}
\usepackage{booktabs}
\usepackage{enumitem}
\usepackage{tcolorbox}
\usepackage{hyperref}
\usepackage{listings}
\usepackage{xcolor}
\usepackage{parskip}
\usepackage{pgfplots}

\usetikzlibrary{arrows.meta}
\pgfplotsset{compat=1.18}

\lstset{
  language=Python,
  basicstyle=\ttfamily\footnotesize,
  breaklines=true,
  columns=flexible,
  keepspaces=true,
  commentstyle=\color{gray}\itshape,
  keywordstyle=\color{blue!70!black}\bfseries,
  stringstyle=\color{green!40!black},
  frame=L,
  xleftmargin=2em,
  showstringspaces=false,
  tabsize=4
}

\title{\textbf{Laboratorio 8: Satisfacion \& Search} \vspace{1em}\\
  Inteligencia Artificial - CC3085}
\author{José Antonio Mérida Castejón}
\date{\today}

\begin{document}

\maketitle

\subsection*{Task 1 - Teoría}
\textit{Responda las siguientes preguntas de forma clara, concisa y argumentada. No basta con definir: se espera que usted
conecte el concepto con sus implicaciones prácticas o estratégicas.}

\begin{enumerate}
  \item \textit{En clase discutimos dos heurísticas vitales: MCV (Variable más Restringida) y LCV (Valor Menos Restringido). Explique con
    sus propias palabras por qué la heurística MCV busca que el algoritmo ``falle lo más rápido posible'', mientras que la heurística
  LCV busca que el algoritmo ``tenga éxito lo más rápido posible''. ¿Por qué esta combinación es tan poderosa en un árbol de Backtracking?}

    \vspace{0.5em}

Uno de los mayores defectos del algoritmo ``clásico'' de \textit{Backtracking} es el tener que 
retroceder luego de recorrer una alta profundidad. La heurística \textit{MCV} nos ayuda a poder 
podar las ramas dónde es poco probable tener éxito y únicamente debemos de explorar un número 
limitado de variables. Por ejemplo, si dentro de un árbol sabemos que una rama 
tiene únicamente 2 caminos posibles tenemos una búsqueda relativamente pequeña, dónde es más 
probable que fallemos en profundidades tempranas en comparación a explorar una rama con 50 caminos posibles.

Sin embargo, encontrar el error rápido no es suficiente por sí solo. La heurística \textit{LCV} 
complementa a \textit{MCV} eligiendo, para cada variable, el valor que menos restringe a las 
variables vecinas. Así, cuando sí existe una solución, el algoritmo tiene más probabilidad de 
encontrarla sin necesidad de retroceder. Mientras \textit{MCV} busca fallar rápido para podar 
ramas inútiles, \textit{LCV} busca tener éxito rápido preservando la mayor cantidad de opciones 
posibles. La combinación de ambas heurísticas es lo que hace al \textit{Backtracking} eficiente 
en la práctica.

  \item \textit{Suponga que está diseñando un CSP para asignar turnos a médicos en un hospital. Tiene una restricción que dice:
    ``El Dr. Pérez y la Dra. Gómez no pueden trabajar en el mismo turno, a menos que sea fin de semana''}

    \begin{itemize}
      \item \textit{¿Cómo modelaría esto utilizando la notación de Factores vista en clase? Defina cuáles serían las variables,
        sus dominios y escriba lógicamente la condición que hace que el factor devuelva 0 o 1}
    \end{itemize}

      \item[] \textbf{Variables y dominios:}
      \begin{itemize}
        \item $X_P \in \{\text{mañana},\ \text{tarde},\ \text{noche},\ \text{libre}\}$: turno asignado al Dr. Pérez.
        \item $X_G \in \{\text{mañana},\ \text{tarde},\ \text{noche},\ \text{libre}\}$: turno asignado a la Dra. Gómez.
        \item $D \in \{\text{semana},\ \text{fin de semana}\}$: indica el tipo de día.
      \end{itemize}

      \vspace{0.5em}
      \textbf{Factor:}
      \[
        f(X_P,\ X_G,\ D) =
        \begin{cases}
          1 & \text{si } D = \text{fin de semana} \\
          1 & \text{si } D = \text{semana} \land X_P \neq X_G \\
          0 & \text{si } D = \text{semana} \land X_P = X_G
        \end{cases}
      \]

      Es decir, el factor devuelve $0$ únicamente cuando es día de semana y ambos médicos
      tienen asignado el mismo turno, violando la restricción.
    \end{itemize}

\end{enumerate}


\end{document}
